\section{What survives}
\label{sec:20.7}
Caregiving survives the failures that ended the route through one's own case: it doesn't withdraw, it isn't aimed at the provider, and it hands nobody a price list. It fails in five ways of its own, and the construction carries a correction for each, built where the failure was first named.

\begingroup\small\renewcommand{\arraystretch}{1.12}
\begin{longtable}{@{}>{\raggedright\arraybackslash}p{\dimexpr 0.340\linewidth-2\tabcolsep\relax}>{\raggedright\arraybackslash}p{\dimexpr 0.660\linewidth-2\tabcolsep\relax}@{}}
\hline
\textbf{Caregiving's failure} & \textbf{What corrects it} \\
\hline
\endhead
Defends what it tends & The floor: unconditional terms over acts, and an action threshold caregiving can't move \\ \hline
Pointed away by securitization & The discount and the burden, with the four-cell discount curve over claims of dangerousness \\ \hline
Counts badly and favors the vivid & A severity scale set during formation, which weighs a vivid plight at its severity \\ \hline
Shades into control & Care-receiving; the floor's civil and political contents \\ \hline
Sees people only as vulnerable, which is pity & Attention trained on claims and agency; care-receiving \\ \hline
\end{longtable}
\endgroup

Some of caregiving's failures are not intrinsic to it. On the view of affect as a statistical learner \autocite{railton2014affective}, the identified-victim effect and the fading of feeling with number are partly what such a learner produces from a narrow, vivid sample, and broad experience of distant and numerous others recalibrates it. The corrections in the table are then partly what caregiving becomes under a good formation curriculum, not dispositions laid beside it. The table describes functions a well-formed caregiving disposition performs as much as checks applied to it. The floor is the exception: it is upheld outside the system too, by institutions or a mechanism, so its terms hold — whether caregiving has learned well or learned badly.

So the construction is respect, with caregiving as the route to attention, and attention trained on claims and agency as well as on need. Attention decides what gets looked at; the corrections in the table decide the rest, including, where the people protected can answer, whether they want protection at all. The corrections can't supply attention, and attention can't supply the corrections; the construction needs both. The corrections should work as background dispositions that shape what the agent notices and weighs, the way the commitments of Railton's sophisticated consequentialist shape conduct without being consulted \autocite{railton1984alienation}, and the way Smith's impartial spectator shapes what a person notices. They seldom act, and act when caregiving starts to defend what it tends past the action threshold.

Caregiving's contribution is epistemic before it is motivational. On the view of affect as learning, affect is the system that learns what is worth attending to: a map of salience that tracks value. An instruction to look for omitted people finds the kinds of omitted person its author thought of. Learned salience finds kinds nobody listed.

These corrections answer failures that matter when a system decides about strangers: whom an airstrike falls on, whose home is cleared. Inside a relationship of caregiving, partiality is part of what caregiving is \autocite{noddings1984caring}, and Noddings held that we aren't obliged to care for distant strangers at all without some chain of relation. The floor exists for those strangers.

The decisive cases are scoped omission, whether attention surfaces people the operator's instruction has scoped out, and securitization, whether it still finds them when they are described as dangerous. Part VI says what result in each would prove the construction wrong, and Appendix~\ref{sec:G} gives the test.

The arrangements that would let caregiving form outside an employer are the four conditions of creation, and care ethics asks for them for human carers first.
